% =====================================================================
%  <TITRE DE L'EXPOSÉ>
%  <occasion> · <public> · <durée> · slides <langue> / oral <langue>
%
%  Plan : plan_presentation.md          Notes : notes_orateur.md
%  Émis par /slides le <date>.          Support : notes_techniques.md
%
%  Compilation : pdflatex deck.tex (×2, pour la table des matières)
%  Contrôle    : slide_audit.py deck.tex --duree <n> --pixels
% =====================================================================
\documentclass[aspectratio=169,11pt]{beamer}
\usetheme[academique]{guyeux}          % ou [affirmee] : audition, keynote, soutenance
\usepackage[french]{babel}             % ou [english]
\usepackage{graphicx}
\usetikzlibrary{calc}
\graphicspath{{figures/}{../../article/figures/}{../../résultats/}}

\title{<Le message, en une phrase nominale>}
\subtitle{<précision, ou rien>}
\author{Christophe Guyeux}
\institute{FEMTO-ST, UMR 6174 CNRS \textperiodcentered\ Universit\'e Marie et Louis Pasteur}
\date{<date de l'exposé>}

\begin{document}

\begin{frame}[plain]\titlepage\end{frame}

% =====================================================================
\section{<Acte 1 : la mise en place>}
% presentation: maillon M1 | budget <n> min | <n> slides

% ---------------------------------------------------------------------
% Slide 2 — <message, qui devient le titre>
% Rôle    : <rôle dans l'arc>
% Durée   : <n> min · Modalité : <forme choisie>
\begin{frame}{<Le message, formulé comme une assertion>}
  \slidefigure{<fichier>}{<ce qu'il faut voir, pas ce que c'est>}
\end{frame}

% =====================================================================
\section{<Acte 2 : la démonstration>}
% presentation: maillons M2,M3 | bascule B1 -> slide de preuve | budget <n> min

% Slide de preuve d'un point de bascule : une figure due, pas une liste.
\begin{frame}{<L'assertion que cette preuve établit>}
  \slidecompare{<état 1>}{%
    <contenu, court>}{<état 2>}{%
    <contenu, court>}
  \kb{<la phrase à retenir, moins de 25 mots>}
\end{frame}

% Un seul nombre porte le message.
\begin{frame}{<L'assertion>}
  \bigchiffre{<chiffre>}{<la phrase qui l'explique>}
\end{frame}

% Le moment où la salle lève les yeux. Deux ou trois par exposé, pas plus.
\begin{sliderupture}
<L'assertion centrale, en une phrase.>
\end{sliderupture}

% =====================================================================
\section{<Acte 3 : la portée>}
% presentation: budget <n> min

\begin{frame}{<Ce que tout cela veut dire>}
  % La slide la plus importante : ne pas répéter les résultats, montrer
  % l'image intégrée.
  \nouveau{<en une ligne>}{<ce qui est neuf par rapport à la littérature>}
\end{frame}

\begin{frame}{<Ce qui reste ouvert>}
  % Réserves et perspectives. Les dire soi-même vaut mieux que se les
  % faire dire. Source : pistes.md.
  \fragilite{<le fait le plus fragile>}{<amplitude, et ce qui le consoliderait>}
\end{frame}

% =====================================================================
\appendix
% Les annexes sont gratuites : elles ne coûtent rien au budget de temps et
% elles font la différence en questions. En préparer généreusement, et les
% indexer dans notes_orateur.md pour les retrouver en direct.

\begin{frame}{Annexe — <sujet>}
\end{frame}

\end{document}
